% Fuente: cuadros 1.6.1 y 1.6.6 de ambos anexos al mensaje presidencial.
% Las cifras originales están en miles de millones; se dividen entre 1.000.
\begin{frame}{PGN 2027: Petro frente a Abelardo}
{\small Proyectos de presupuesto. Billones de pesos corrientes.}
\vspace{0.12cm}
\begin{center}
\begingroup
\fontsize{10.2}{12.2}\selectfont
\setlength{\tabcolsep}{7pt}
\renewcommand{\arraystretch}{1.12}
\begin{tabular}{@{}lrrr@{}}
\toprule
 & \textbf{Petro} & \textbf{Abelardo} & \textbf{Diferencia} \\
 & julio de 2026 & agosto de 2026 & Abelardo $-$ Petro \\
\midrule
Funcionamiento & 367,698 & 392,560 & $+24{,}862$ \\
Inversión & 89,960 & 86,960 & $-3{,}000$ \\
\addlinespace[3pt]
\textbf{Servicio de la deuda} & \textbf{118,041} & \textbf{155,432} & $\mathbf{+37{,}391}$ \\
\quad Amortizaciones (principal) & 27,604 & 58,272 & $+30{,}668$ \\
\quad Intereses & 87,714 & 94,421 & $+6{,}707$ \\
\quad Comisiones y otros gastos & 0,295 & 0,311 & $+0{,}016$ \\
\quad Fondo de Contingencias & 2,429 & 2,429 & $0{,}000$ \\
\midrule
\textbf{Total PGN} & \textbf{575,699} & \textbf{634,952} & $\mathbf{+59{,}253}$ \\
\bottomrule
\end{tabular}
\endgroup
\end{center}
\vspace{-0.06cm}
{\fontsize{7.2}{8.6}\selectfont
Fuente: MHCP--DGPPN, anexos al mensaje presidencial de julio y agosto de 2026,
cuadros 1.6.1 y 1.6.6. Cifras proyectadas, no ejecución.\par
El principal incluye Acuerdos Marco de Retribución. Se conservan los totales publicados:
la suma de los componentes de deuda los excede en 0,001 billones en ambos proyectos.\par}
\note{[Sources] MHCP--DGPPN, Anexo al Mensaje del PGN 2027, edición julio de 2026,
cuadro 1.6.1, página impresa 37 (PDF 58), y cuadro 1.6.6, página impresa 53 (PDF 74).
Archivo fuente: C:/Users/olive/Downloads/2.Anexo Mensaje PGN 2027.pdf.
MHCP--DGPPN, Anexo al Mensaje del PGN 2027 reformulado, edición agosto de 2026,
cuadro 1.6.1, página impresa 45 (PDF 58), y cuadro 1.6.6, página impresa 61 (PDF 74).
Archivo fuente: C:/Users/olive/Downloads/2. Anexo al Mensaje.pdf.
Ambas columnas corresponden a proyectos del PGN para la vigencia 2027, no a gasto
fiscal del GNC ni a ejecución. Se divide la unidad original entre 1.000. La diferencia
se calcula sobre los valores publicados, sin ajustar sus discrepancias de suma.
El título del cuadro 1.6.6 dice 2025--2026 en ambas fuentes, pero sus columnas
identifican 2026 y Proyecto 2027: se usa exclusivamente la columna Proyecto 2027.
El principal incluye Acuerdos Marco de Retribución por 0,490 y 0,990 billones,
respectivamente. Las cifras se cotejaron con las tablas; no con el párrafo bajo
el cuadro de agosto, que agrupa incorrectamente intereses y otros rubros.
Consulta y cotejo: 2 de septiembre de 2026.}
\end{frame}
